\documentclass{amsart}
% For dealing with references we use the comment environment
\usepackage{verbatim}
\newenvironment{reference}{\comment}{\endcomment}
%\newenvironment{reference}{}{}
\newenvironment{slogan}{\comment}{\endcomment}
\newenvironment{history}{\comment}{\endcomment}

% For commutative diagrams we use Xy-pic
\usepackage[all]{xy}

% We use 2cell for 2-commutative diagrams.
\xyoption{2cell}
\UseAllTwocells

% We use multicol for the list of chapters between chapters
\usepackage{multicol}

% This is generally recommended for better output
\usepackage{lmodern}
\usepackage[T1]{fontenc}

% For cross-file-references

% Package for hypertext links:
\usepackage{hyperref}

% For any local file, say "hello.tex" you want to link to please

% Theorem environments.
%
\theoremstyle{plain}
\newtheorem{theorem}[subsection]{Theorem}
\newtheorem{proposition}[subsection]{Proposition}
\newtheorem{lemma}[subsection]{Lemma}

\theoremstyle{definition}
\newtheorem{definition}[subsection]{Definition}
\newtheorem{example}[subsection]{Example}
\newtheorem{exercise}[subsection]{Exercise}
\newtheorem{situation}[subsection]{Situation}

\theoremstyle{remark}
\newtheorem{remark}[subsection]{Remark}
\newtheorem{remarks}[subsection]{Remarks}

\numberwithin{equation}{subsection}

% Macros
%
\def\lim{\mathop{\mathrm{lim}}\nolimits}
\def\colim{\mathop{\mathrm{colim}}\nolimits}
\def\Spec{\mathop{\mathrm{Spec}}}
\def\Hom{\mathop{\mathrm{Hom}}\nolimits}
\def\Ext{\mathop{\mathrm{Ext}}\nolimits}
\def\SheafHom{\mathop{\mathcal{H}\!\mathit{om}}\nolimits}
\def\SheafExt{\mathop{\mathcal{E}\!\mathit{xt}}\nolimits}
\def\Sch{\mathit{Sch}}
\def\Mor{\mathop{\mathrm{Mor}}\nolimits}
\def\Ob{\mathop{\mathrm{Ob}}\nolimits}
\def\Sh{\mathop{\mathit{Sh}}\nolimits}
\def\NL{\mathop{N\!L}\nolimits}
\def\CH{\mathop{\mathrm{CH}}\nolimits}
\def\proetale{{pro\text{-}\acute{e}tale}}
\def\etale{{\acute{e}tale}}
\def\QCoh{\mathit{QCoh}}
\def\Ker{\mathop{\mathrm{Ker}}}
\def\Im{\mathop{\mathrm{Im}}}
\def\Coker{\mathop{\mathrm{Coker}}}
\def\Coim{\mathop{\mathrm{Coim}}}

% Boxtimes
%
\DeclareMathSymbol{\boxtimes}{\mathbin}{AMSa}{"02}

%
% Macros for moduli stacks/spaces
%
\def\QCohstack{\mathcal{QC}\!\mathit{oh}}
\def\Cohstack{\mathcal{C}\!\mathit{oh}}
\def\Spacesstack{\mathcal{S}\!\mathit{paces}}
\def\Quotfunctor{\mathrm{Quot}}
\def\Hilbfunctor{\mathrm{Hilb}}
\def\Curvesstack{\mathcal{C}\!\mathit{urves}}
\def\Polarizedstack{\mathcal{P}\!\mathit{olarized}}
\def\Complexesstack{\mathcal{C}\!\mathit{omplexes}}
% \Pic is the operator that assigns to X its picard group, usage \Pic(X)
% \Picardstack_{X/B} denotes the Picard stack of X over B
% \Picardfunctor_{X/B} denotes the Picard functor of X over B
\def\Pic{\mathop{\mathrm{Pic}}\nolimits}
\def\Picardstack{\mathcal{P}\!\mathit{ic}}
\def\Picardfunctor{\mathrm{Pic}}
\def\Deformationcategory{\mathcal{D}\!\mathit{ef}}

\setlength{\emergencystretch}{3em}
\begin{document}
\title[Stacks Project, Tag 0D2I]{305 --- Patching equivalence for flat modules over a fibre product}
\maketitle
\noindent Copyright (C) 2005--2025 Johan de Jong. Stacks Project authors. Source: \href{https://stacks.math.columbia.edu/tag/0D2I}{Tag 0D2I}.

\noindent Editorial difficulty note (not author text). Difficulty H2: The central obligation is to prove flatness of the patched fibre-product module for arbitrary ideals. The author reduces through the kernel ideal, proves Tor vanishing, saturates an ideal as an A-prime-module and recovers injectivity; the equivalence follows only after this substantive tensor argument..

\noindent Editorial provenance note (not author text). History: excerpt from revision \texttt{a0444\allowbreak 6e57e\allowbreak c1fbc\allowbreak 252a8\allowbreak 71afc\allowbreak ec775\allowbreak 2fb28\allowbreak 07b14}, more-algebra.tex, lines 806--889. Reference presentation was adapted; original-author context and core support blocks were retained or added. The mathematical wording is unchanged. Licensed under GNU FDL 1.2 or later; no invariant sections or cover texts. See ../licenses/COPYING. Context blocks reproduce author definitions and supporting results; they are not additional counted problems.

\begin{situation}
\label{situation-module-over-fibre-product}
In the following we will consider ring maps
$$
\xymatrix{
B \ar[r] & A & A' \ar[l]
}
$$
where we assume $A' \to A$ is surjective with kernel $I$.
In this situation we set $B' = B \times_A A'$ to
obtain a cartesian square
$$
\xymatrix{
A & A' \ar[l] \\
B \ar[u] & B' \ar[l] \ar[u]
}
$$
\end{situation}

\begin{lemma}
\label{lemma-flat-module-over-fibre-product}
With $A, A', B, B', I$ as in
Situation \ref{situation-module-over-fibre-product}.
\begin{enumerate}
\item Let $(N, M', \varphi)$ be an object of
$\text{Mod}_B \times_{\text{Mod}_A} \text{Mod}_{A'}$.
If $M'$ is flat over $A'$ and $N$ is flat over $B$, then
$N' = N \times_{\varphi, M} M'$ is flat over $B'$.
\item If $L'$ is a flat $B'$-module, then
$L' = (L \otimes_{B'} B) \times_{(L \otimes_{B'} A)} (L \otimes_{B'} A')$.
\item The category of flat $B'$-modules is equivalent to the
full subcategory of $\text{Mod}_B \times_{\text{Mod}_A} \text{Mod}_{A'}$
consisting of triples
$(N, M', \varphi)$ with $N$ flat over $B$ and $M'$ flat over $A'$.
\end{enumerate}
\end{lemma}

\begin{proof}
In the proof we will use Lemma \href{https://stacks.math.columbia.edu/tag/07RU}{lemma module over fibre product (Tag 07RU)}
without further mention.

\medskip\noindent
Proof of (1). Set $J = \Ker(B' \to B)$. This is an ideal of $B'$
mapping isomorphically to $I = \Ker(A' \to A)$.
Let $\mathfrak b' \subset B'$ be an ideal.
We have to show that $\mathfrak b' \otimes_{B'} N' \to N'$
is injective, see Algebra, Lemma \href{https://stacks.math.columbia.edu/tag/00HD}{lemma flat (Tag 00HD)}.
We know that
$$
\mathfrak b'/(\mathfrak b' \cap J) \otimes_{B'} N' =
\mathfrak b'/(\mathfrak b' \cap J) \otimes_B N \to N
$$
is injective as $N$ is flat over $B$. As
$\mathfrak b' \cap J \to \mathfrak b' \to
\mathfrak b'/(\mathfrak b' \cap J) \to 0$ is exact, we
conclude that it suffices to show that
$(\mathfrak b' \cap J) \otimes_{B'} N' \to N'$
is injective. Thus we may assume that $\mathfrak b' \subset J$.
Next, since $J \to I$ is an isomorphism we have
$$
J \otimes_{B'} N' =
I \otimes_{A'} A' \otimes_{B'} N' =
I \otimes_{A'} M'
$$
which maps injectively into $M'$ as $M'$ is a flat $A'$-module.
Hence $J \otimes_{B'} N' \to N'$ is injective and we conclude
that $\text{Tor}_1^{B'}(B'/J, N') = 0$, see
Algebra, Remark \href{https://stacks.math.columbia.edu/tag/00M6}{remark Tor ring mod ideal (Tag 00M6)}.
Thus we may apply Algebra, Lemma \href{https://stacks.math.columbia.edu/tag/051C}{lemma what does it mean (Tag 051C)}
to $N'$ over $B'$ and the ideal $J$.
Going back to our ideal $\mathfrak b' \subset J$, let
$\mathfrak b' \subset \mathfrak b'' \subset J$ be the smallest ideal
whose image in $I$ is an $A'$-submodule of $I$. In other words,
we have $\mathfrak b'' = A' \mathfrak b'$ if we view $J = I$
as $A'$-module. Then $\mathfrak b''/\mathfrak b'$ is killed
by $J$ and we get a short exact sequence
$$
0 \to \mathfrak b' \otimes_{B'} N' \to
\mathfrak b'' \otimes_{B'} N' \to
\mathfrak b''/\mathfrak b' \otimes_{B'} N' \to 0
$$
by the vanishing of $\text{Tor}_1^{B'}(\mathfrak b''/\mathfrak b', N')$
we get from the application of the lemma.
Thus we may replace $\mathfrak b'$ by $\mathfrak b''$.
In particular we may assume $\mathfrak b'$ is an $A'$-module
and maps to an ideal of $A'$. Then
$$
\mathfrak b' \otimes_{B'} N' =
\mathfrak b' \otimes_{A'} A' \otimes_{B'} N' =
\mathfrak b' \otimes_{A'} M'
$$
This tensor product maps injectively into $M'$ by our assumption that
$M'$ is flat over $A'$. We conclude that
$\mathfrak b' \otimes_{B'} N' \to N' \to M'$ is injective
and hence the first map is injective as desired.

\medskip\noindent
Proof of (2). This follows by tensoring the short exact sequence
$0 \to B' \to B \oplus A' \to A \to 0$ with $L'$ over $B'$.

\medskip\noindent
Proof of (3). Immediate consequence of (1) and (2).
\end{proof}
\end{document}
